\begin{apartats}

\apartat{1,25}
Determina els valors de $a$ i $b$ perquè la funció següent sigui contínua a tots els
punts de $\mathbb{R}$.
\[
f(x)=\begin{cases}
  x^2-a & \si{x<-1},\\[3pt]
  bx+2  & \si{-1\le x\le 2},\\[3pt]
  ax+b  & \si{x>2}.
\end{cases}
\]

\begin{solucio}
Cada branca és contínua al seu tros, de manera que només cal imposar la continuïtat
als enganxaments.\\
En $x=-1$: $\lim_{x\to-1^-}f(x)=1-a$ i $f(-1)=-b+2$, d'on $1-a=-b+2$, és a dir $b=a+1$.\\
En $x=2$: $f(2)=2b+2$ i $\lim_{x\to2^+}f(x)=2a+b$, d'on $2b+2=2a+b$, és a dir $b=2a-2$.\\
Igualant: $a+1=2a-2\Rightarrow \boxed{a=3}$ i $\boxed{b=4}$.
\end{solucio}

\apartat{1,25}
Estudia la continuïtat de la funció $h$ en $x=0$, $x=2$ i $x=3$, i classifica'n els
tipus de discontinuïtat que presenta.
\[
h(x)=\begin{cases}
  2x+1 & \si{x<0},\\[3pt]
  e^{x} & \si{0\le x\le 2},\\[3pt]
  \dfrac{x-2}{x^2-5x+6} & \si{x>2}.
\end{cases}
\]

\begin{solucio}
Per a $x>2$, $\dfrac{x-2}{(x-2)(x-3)}=\dfrac{1}{x-3}$.\\
$x=0$: laterals $1$ i $e^0=1$, i $h(0)=1$. \textbf{És contínua.}\\
$x=2$: $h(2)=e^2$ i $\lim_{x\to2^-}h(x)=e^2$, però $\lim_{x\to2^+}h(x)=\frac{1}{2-3}=-1$.
Laterals finits i diferents: \textbf{salt finit}.\\
$x=3$: $h(3)$ no existeix i els laterals valen $-\infty$ i $+\infty$:
\textbf{salt infinit} (asímptota vertical $x=3$).
\end{solucio}

\end{apartats}
